% ID: SEL-D0041-OB-E02
% ESTADO: validada. S042, cotejada con examen y criterios.
\textbf{a)} Los dos tramos son continuos en su interior.
En $2$, el primero da $-4+2+6=4$ y el límite derecho
es $2+2=4$. Por tanto $f$ es continua en todo $\mathbb R$.

\textbf{b)} Las pendientes de los tramos son $1-2x$ y 1.
En $2$ las pendientes laterales son $-3$ y 1. Así que
$f$ es derivable en $\mathbb R\setminus\{2\}$.

\textbf{c)} La derivada del primer tramo se anula en $1/2$
y pasa de positiva a negativa. Por ello $f$ crece en
$(-\infty,1/2]$, decrece en $[1/2,2]$ y vuelve a crecer
en $[2,\infty)$. Tiene máximo relativo en $(1/2,25/4)$ y
mínimo relativo en $(2,4)$. Los dos extremos son locales,
pues la función no está acotada ni superior ni inferiormente.
\[
-x^2+x+6=-(x-3)(x+2).
\]
Sólo $-2$ pertenece al tramo cuadrático. El cero de la recta,
$-2$, está fuera del tramo lineal. Los cortes son
$\boxed{(-2,0)}$ y $\boxed{(0,6)}$.
\begin{center}
\setlength{\unitlength}{7mm}
\begin{picture}(12,9)(-4,-4)
\qbezier(-3.5,0)(1,0)(5.5,0)\put(5.6,0){$x$}
\qbezier(0,-3.5)(0,0.5)(0,4.5)\put(0,4.6){$y$}
\qbezier(-3,-3)(-0.5,5.75)(2,2)
\qbezier(2,2)(3.5,2.75)(5,3.5)
\put(0.5,3.35){\makebox(0,0)[b]{$(1/2,25/4)$}}
\put(2.1,1.7){$(2,4)$}
\put(-2,-0.3){\makebox(0,0)[t]{$(-2,0)$}}
\end{picture}
\end{center}
